\section{USAMO 2014}
        \begin{enumerate}
        	\item (\textit{USAMO 2014}) Let $a,b,c,d$ be real numbers such that $b-d \ge 5$ and all zeros $x_1, x_2, x_3,$ and $x_4$ of the polynomial $P(x)=x^4+ax^3+bx^2+cx+d$ are real. Find the smallest value the product 
 \[(x_1^2+1)(x_2^2+1)(x_3^2+1)(x_4^2+1)\] can take.


 



\textbf{Solution}

Using the hint we turn the equation into $\prod_{k=1} ^4 (x_k-i)(x_k+i) \implies P(i)P(-i) \implies (b-d-1)^2 + (a-c)^2 \implies \boxed{16}$. This minimum is achieved when all the $x_i$ are equal to $1$.

 A more detailed version goes as follows:

 Observe that $x^2+1=x^2-(-1)=x^2-i^2=(x-i)(x+i).$ Now, notice that:
 \[\prod_{k=1} ^4 (x_k-i)(x_k+i)=\prod_{k=1}^4(x_k-i)\cdot\prod_{k=1}^4(x_k+i).\]The definition of $P(x)$ is $(x-x_1)(x-x_2)(x-x_3)(x-x_4)$ where the leading coeffecient is $1$ since $P$ is a monic polynomial as given in the problem. Then, substituting in $i,$ we have:
 \[P(i)=(i-x_1)\cdots(i-x_4)=(-1)^4(x_1-i)\cdots(x_4-i).\]This is exactly our first product. Our second product can be found as follows:
 \[P(-i)=(-i-x_1)\cdots(-i-x_4)=(-1)^4(x_1+i)\cdots(x_4+i).\]Hence, what we want to find is $P(i)P(-i).$ But substituting in $i$ into our expressoin $x^4+ax^3+bx^2+cx+d$ gets us $P(i)=(1-b+d)+(c-a)i,$ and $P(-i)=(1-b+d)+(a-c)i.$ Hence:
 \[P(i)P(-i)=((1-b+d)+(c-a)i)((1-b+d)+(a-c)i)=(1-b+d)^2-(a-c)^2=(-1)^2(b-d-1)^2-(a-c)^2 \ge 4^2 -0^2=16\]where equality holds at $a=c$ and $b-d=5 \Rightarrow b=d+5.$ Hence, the minimum is $\boxed{16}.$ 

 To finish this off we find a construction for this minimum. We know that $a=c$ and $b=d+5.$ Hence
 \[P(x)=x^4+cx^3+(d+5)x^2+cx+d.\]We set $d=1$ to get as much symmetry as possible within our polynomial. This leads to $x^4+cx^3+6x^2+cx+1.$ However, note that $\binom{4}{2}=6$ and $\binom{4}{0}=\binom{4}{4}=1,$ so with some wishful thinking this leads us to think about the binomial theorem. We can try $(x\pm 1)^4$ and realize that those solutions do work. 

 Hence, $16$ is obtained when all of $x_i$ are equal to $1$ or all equal to $-1.$

 


	\item (\textit{USAMO 2014}) Let $\mathbb{Z}$ be the set of integers. Find all functions $f : \mathbb{Z} \rightarrow \mathbb{Z}$ such that 
 \[xf(2f(y)-x)+y^2f(2x-f(y))=\frac{(f(x))^2}{x}+f(yf(y))\] for all $x, y \in \mathbb{Z}$ with $x \neq 0$.


 



\textbf{Solution}

Note: This solution is kind of rough. I didn't want to put my 7-page solution all over again. It would be nice if someone could edit in the details of the expansions.

 Lemma 1: $f(0) = 0$.Proof: Assume the opposite for a contradiction. Plug in $x = 2f(0)$ (because we assumed that $f(0) \neq 0$), $y = 0$. What you get eventually reduces to:
 \[4f(0)-2 = \left( \frac{f(2f(0))}{f(0)} \right)^2\]which is a contradiction since the LHS is divisible by 2 but not 4.

 Then plug in $y = 0$ into the original equation and simplify by Lemma 1. We get:
 \[x^2f(-x) = f(x)^2\]Then:

 
 \begin{align*} x^6f(x) &= x^4\bigl(x^2f(x)\bigr)\\ &= x^4\bigl((-x)^2f(-(-x))\bigr)\\ &= x^4(-x)^2f(-(-x))\\ &= x^4f(-x)^2\\ &= f(x)^4 \end{align*}


 Therefore, $f(x)$ must be 0 or $x^2$.

 Now either $f(x)$ is $x^2$ for all $x$ or there exists $a \neq 0$ such that $f(a)=0$. The first case gives a valid solution. In the second case, we let $y = a$ in the original equation and simplify to get:
 \[xf(-x) + a^2f(2x) = \frac{f(x)^2}{x}\]But we know that $xf(-x) = \frac{f(x)^2}{x}$, so:
 \[a^2f(2x) = 0\]Since $a$ is not 0, $f(2x)$ is 0 for all $x$ (including 0). Now either $f(x)$ is 0 for all $x$, or there exists some $m \neq 0$ such that $f(m) = m^2$. Then $m$ must be odd. We can let $x = 2k$ in the original equation, and since $f(2x)$ is 0 for all $x$, stuff cancels and we get:
 \[y^2f(4k - f(y)) = f(yf(y))\]\textbf{for $\mathbf{k \neq 0}$.}Now, let $y = m$ and we get:
 \[m^2f(4k - m^2) = f(m^3)\]Now, either both sides are 0 or both are equal to $m^6$. If both are $m^6$ then:
 \[m^2(4k - m^2)^2 = m^6\]which simplifies to:
 \[4k - m^2 = \pm m^2\]Since $k \neq 0$ and $m$ is odd, both cases are impossible, so we must have:
 \[m^2f(4k - m^2) = f(m^3) = 0\]Then we can let $k$ be anything except 0, and get $f(x)$ is 0 for all $x \equiv 3 \pmod{4}$ except $-m^2$.  Also since $x^2f(-x) = f(x)^2$, we have $f(x) = 0 \Rightarrow f(-x) = 0$, so $f(x)$ is 0 for all $x \equiv 1 \pmod{4}$ except $m^2$. So $f(x)$ is 0 for all $x$ except $\pm m^2$. Since $f(m) \neq 0$, $m = \pm m^2$. Squaring, $m^2 = m^4$ and dividing by $m$, $m = m^3$. Since $f(m^3) = 0$, $f(m) = 0$, which is a contradiction for $m \neq 1$. However, if we plug in $x = 1$ with $f(1) = 1$ and $y$ as an arbitrary large number with $f(y) = 0$ into the original equation, we get $0 = 1$ which is a clear contradiction, so our only solutions are $f(x) = 0$ and $f(x) = x^2$.

 



\textbf{Solution 2}

Given that the range of f consists entirely of integers, it is clear that the LHS must be an integer and $f(yf(y))$ must also be an integer, therefore $\frac{f(x)^2}{x}$ is an integer. If $x$ divides $f(x)^2$ for all integers $x \ne 0$, then $x$ must be a factor of $f(x)$, therefore $f(0)=0$. Now, by setting $y=0$ in the original equation, this simplifies to $xf(-x)=\frac{f(x)^2}{x}$. Assuming $x \ne 0$, we have $x^2f(-x)=f(x)^2$. Substituting in $-x$ for $x$ gives us $x^2f(x)=f(-x)^2$. Substituting in $\frac{f(x)^2}{x^2}$ in for $f(-x)$ in the second equation gives us $x^2f(x)=\frac{f(x)^4}{x^4}$, so $x^6f(x)=f(x)^4$. In particular, if $f(x) \ne 0$, then we have $f(x)^3=x^6$, therefore $f(x)$ is equivalent to $0$ or $x^2$ for every integer $x$. Now, we shall prove that if for some integer $t \ne 0$, if $f(t)=0$, then $f(x)=0$ for all integers $x$. If we assume $f(y)=0$ and $y \ne 0$ in the original equation, this simplifies to $xf(-x)+y^2f(2x)=\frac{f(x)^2}{x}$. However, since $x^2f(-x)=f(x)^2$, we can rewrite this equation as $\frac{f(x)^2}{x}+y^2f(2x)=\frac{f(x)^2}{x}$, $y^2f(2x)$ must therefore be equivalent to $0$. Since, by our initial assumption, $y \ne 0$, this means that $f(2x)=0$, so, if for some integer $y \ne 0$, $f(y)=0$, then $f(x)=0$ for all integers $x$. The contrapositive must also be true, i.e. If $f(x) \ne 0$ for all integers $x$, then there is no integral value of $y \ne 0$ such that $f(y)=0$, therefore $f(x)$ must be equivalent for $x^2$ for every integer $x$, including $0$, since $f(0)=0$. Thus, $f(x)=0, x^2$ are the only possible solutions.

 



\textbf{Solution 3}

Let's assume $f(0)\neq 0.$ Substitute $(x,y)=(2f(0),0)$ to get 
 \[2f(0)^2=f(2f(0))^2/2f(0)+f(0)\]
 \[2f(0)^2(2f(0)-1)=f(2f(0))^2\]

 This means that $2(2f(0)-1)$ is a perfect square. However, this is impossible, as it is equivalent to $2\pmod{4}.$ Therefore, $f(0)=0.$ Now substitute $x\neq 0, y=0$ to get 
 \[xf(-x)=\frac{f(x)^2}{x} \implies x^2f(-x)=f(x)^2.\]Similarly, 
 \[x^2f(x)=f(-x)^2.\]From these two equations, we can find either $f(x)=f(-x)=0,$ or $f(x)=f(-x)=x^2.$ Both of these are valid solutions on their own, so let's see if there are any solutions combining the two. 

 Let's say we can find $f(x)=x^2, f(y)=0,$ and $x,y\neq 0.$ Then 
 \[xf(-x)+y^2f(2x)=f(x)^2/x.\]
 \[y^2f(2x)=x-x^3.\] (NEEDS FIXING: $f(x)^2/x= x^4/x = x^3$, so the RHS is $0$ instead of $x-x^3$.)

 If $f(2x)=4x^2,$ then $y^2=\frac{x-x^3}{4x^2}=\frac{1-x^2}{4x},$ which is only possible when $y=0.$ This contradicts our assumption. Therefore, $f(2x)=0.$ This forces $x=\pm 1$ due to the right side of the equation. Let's consider the possibility $f(2)=0, f(1)=1.$ Substituting $(x,y)=(2,1)$ into the original equation yields  
 \[0=2f(0)+1f(2)=0+f(1)=1,\] which is impossible. So $f(2)=f(-2)=4$ and there are no solutions "combining" $f(x)=x^2$ and $f(x)=0.$

 Therefore our only solutions are $\boxed{f(x)=0}$ and $\boxed{f(x)=x^2.}$

 



\textbf{Solution 4}

Let the given assertion be $P(x, y)$. We try $P(x, 0)$ and get $xf(2c-x)=f(x)^2/x+c$, where $f(0)=c$. We plug in $x=c$ and get $cf(c)=f(c)^2/c+c$. Rearranging and solving for $c^2$ gives us $c^2=\frac{f(c)^2}{f(c)-1}$. Obviously, the only $c$ that works such that the RHS is an integer is $c=0$, and thus $f(0)=0$.

 We use this information on assertion $P(x,0)$ and obtain $xf(-x)=f(x^2)/x$, or $f(-x)=\frac{f(x)^2}{x^2}$. Thus, $f(x)$ is an even function. It follows that $f(x)=0, x^2$ for each $x$. We now prove that $f(x)=x^2$, f(x)=0$ are the only solutions. [in progress] ~SigmaPiE

 


	\item (\textit{USAMO 2014}) Prove that there exists an infinite set of points 
 \[\ldots,\,\,\,\,P_{-3},\,\,\,\,P_{-2},\,\,\,\,P_{-1},\,\,\,\,P_0,\,\,\,\,P_1,\,\,\,\,P_2,\,\,\,\,P_3,\,\,\,\,\ldots\] in the plane with the following property: For any three distinct integers $a,b,$ and $c$, points $P_a$, $P_b$, and $P_c$ are collinear if and only if $a+b+c=2014$.


 



\textbf{Solution (Group Theory)}

Consider an elliptic curve with a generator $g$, such that $g$ is not a root of $0$. By repeatedly adding $g$ to itself under the standard group operation, with can build $g, 2g, 3g, \ldots$ as well as $-g, -2g, -3g, \ldots$. If we let 
 \[P_k = (3k-2014)g\] then we can observe that collinearity between $P_a$, $P_b$, and $P_c$ occurs only if $P_a + P_b + P_c = 0$ (by definition of the group operation), which is equivalent to $(3a-2014)g + (3b-2014)g + (3c-2014)g = (3a+3b+3c-3*2014)g = 0$, or $3a + 3b + 3c = 3*2014$, or $a + b + c = 2014$. We know that all these points $P_k$ exist because $3k-2014$ is never 0 for integer $k$, so that none of these points need to be point at infinity (the identity element of the group).

 



\textbf{Solution 2 (Function Theory)}

Consider letting $P_x$ be the point $(x, f(x))$, where $f(x) = x^3 - 2014x^2$. Then if three points $P_a, P_b, P_c$ are on the same line $y = mx + p$, they must be the solutions to the equation $x^3 - 2014x^2 = mx + p$ (i.e. the intersection of $f$ and the line). By Vieta's Formulas, the sum of the roots to this equation is 2014. Because each value of x corresponds to a different one of a, b, and c by definition of $P_x$, $a + b + c = 2014$. Conversely, if $a + b + c = 2014$, they must be the solutions to $(x-a)(x-b)(x-c) = x^3 - 2014x^2 - mx - p = 0$ for some real $m$ and $p$. Clearly, then, $P_a, P_b, P_c$ must all lie on the line $y = mx + p$. Hence, our setting $P_x = f(x)$ produces a valid infinite set of points.

 \textit{Note: We could have let $f(x) = ax^3 - 2014ax^2 + bx + c$, where a, b, and c are arbitrary constants. (a is nonzero.)}

 



\textbf{Solution 3}

Consider $a \mapsto a-671$, $b \mapsto b-671$, $c \mapsto c-671$, so we find all cases that satisfy $a+b+c = 1$. We prove infnitely many points satisfy $P(n) = (n, Q(n))$ for a polynomial $Q(x)$. Then, notice that if $(a,Q(a)), (b,Q(b)), (c,Q(c))$ are collinear, then
 \[(b-c)Q(a) + (c-a)Q(b) + (a-b)Q(c) = 0\].

 \textbf{Lemma 1:} $(b-c)Q(a) + (c-a)Q(b) + (a-b)Q(c)$ is divisible by $(1-a-b-c)(a-b)(b-c)(c-a)$.

 \textbf{Proof of Lemma 1:}

 This is trivial by the fact whenever $a+b+c = 1$ the equality to 0 holds. $\blacksquare$

 Indeed, we rewrite as

 
 \[(1-a-b-c)(a-b)(b-c)(c-a)R(a,b,c)\]

 for a multivariable function $R(a,b,c)$.

 \textbf{Lemma 2:} We can assume $R(a,b,c) = 1$.

 \textbf{Proof of Lemma 2:}

 We already can say that $R(x) = k$ which is a constant and now we have $kO(x)$ where $O(x)$ is a polynomial, hence if $O(x) = 0$ whenever $a+b+c = 1$ then one can see that $O(x) = kO(x) = 0$ regardless, hence assume $R(x) = 1 \quad \blacksquare$.

 Therefore,

 
 \[(b-c)Q(a) + (c-a)Q(b) + (a-b)Q(c)=(1-a-b-c)(a-b)(b-c)(c-a)\]. 

 \textbf{2014 USAMO P3:}

 Here, $a^3$ occurs on the RHS and $(b-c)a^3$ is the term. These are equivalent, hence solutions exist only if $Q(x)$ is monic. Then, let $Q(x) = x^3 + R(x)$ where $R(x)$ is a polynomial of degree at most 2. Hence we have that, after reduction,

 
 \[(b-c)R(a) + (c-a)R(b) + (a-b)R(c) = (a-b)(b-c)(c-a)\].

 Comparing leading coefficients we see that $R(x)$ has leading coefficient -1 and is quadratic. It is easy to see $R(x) = -x^2$ works, hence $Q(x) = x^3-x^2$, and there exist infinitely many $P(x)$ that satisfy the coordinates:
 \[P(n) = (n+671, n^3-n^2) \qquad \square\]

 


	\item (\textit{USAMO 2014}) Let $k$ be a positive integer. Two players $A$ and $B$ play a game on an infinite grid of regular hexagons. Initially all the grid cells are empty. Then the players alternately take turns with $A$ moving first. In his move, $A$ may choose two adjacent hexagons in the grid which are empty and place a counter in both of them. In his move, $B$ may choose any counter on the board and remove it. If at any time there are $k$ consecutive grid cells in a line all of which contain a counter, $A$ wins. Find the minimum value of $k$ for which $A$ cannot win in a finite number of moves, or prove that no such minimum value exists.


 



\textbf{Solution}

We claim that the minimum $k$ such that A cannot create a $k$ in a row is $\boxed{6}$.

 It is easy to verify that player A can create a 5 in a row.

 Let A place two counters anywhere, and B take off one of them. Then, A should create a "triangle of hexagons" by placing two adjacent counters also next to the unremoved one, so that no matter what B does there will be a line of two counters on adjacent hexagons.

 \textit{Note: A diagram and more conciseness is required in the following solution. Better maneuvers are always appreciated.}

 Then, it is advisable for A to make a 4-in-a-row by placing 2 counters at the end of the line. B to prevent a 5-in-a-row must counter by removing a middle counter (a counter not on the edge of the 4-in-a-row). Player A then could reinstate the threat by placing a counter there and one adjacent to that hexagon (but not adjacent to the other middle counter). It is not in B's best interest to then remove the other middle counter, for then A can "add" a counter to both hexagons (adjacent to the removed counter) on the same side of the originally placed hexagon while preserving the threat. Now, we have two different threats B cannot both counter. Thus, assume B does not remove the other middle counter. Eventually, by "adding" counters and preserving the threat player A can create a web of hexagons surrounding the chosen middle counter. As we have proven, it will be disastrous for B to then remove the other middle counter; thus, B has to remove the chosen middle counter.

 Now, player A can administer the decisive winning maneuver as follows. Let X be two adjacent counters of the web, one of which is adjacent to the initial edge counter. Let Y be the other two counters that satisfy the condition. The line X should not have more than two counters on it. Player A should place his two counters on the first two hexagons not adjacent to any counter in X (or the web), so that B is forced to remove the outermost counter A played. Then, A should place two counters adjacent to his remaining counter in a direction parallel to line Y. B has to expunge the outermost one; otherwise A can place two counters to complete a 5-in-a-row. Let A place two more counters in the same direction; then, B has to remove the middle counter played. Now, turn in a direction parallel to line X; let A place a counter on the first two spaces adjacent to counters already placed. B will remove the second counter played. A's knockout blow occurs when he reinstates the lost counter and a counter adjacent to that counter and the remaining one he played last turn, and it is easy to verify he is guaranteed a 5-in-a-row. Therefore, A can (after an elongated siege) obtain a 5-in-a-row.

 To prove that 6-in-a-row is impossible, tile the hexagon board in an A-B-C-A-B-C-A-B-C format. Define an "A tile" to be a tile engraved with the letter A. To prevent player A from obtaining a 6 in a row, all player B has to do is to remove counters placed on A tiles by player A. Because no two A tiles are adjacent, player A can play at most one counter on an A tile at a time; thus, there can at most be one counter on A tiles any time of the game. Clearly, then, a 6 in a row, which requires A-B-C-A-B-C (or two counters on A tiles), is impossible, completing the proof.

 


	\item (\textit{USAMO 2014}) Let $ABC$ be a triangle with orthocenter $H$ and let $P$ be the second intersection of the circumcircle of triangle $AHC$ with the internal bisector of the angle $\angle BAC$.  Let $X$ be the circumcenter of triangle $APB$ and $Y$ the orthocenter of triangle $APC$.  Prove that the length of segment $XY$ is equal to the circumradius of triangle $ABC$.


 



\textbf{Solution 1}

Let $O_1$ be the center of $(AHPC)$, $O$ be the center of $(ABC)$. Note that $(O_1)$ is the reflection of $(O)$ across $AC$, so $AO=AO_1$. Additionally
 \[\angle AYC=180-\angle APC=180-\angle AHC=\angle B\]so $Y$ lies on $(O)$. Now since $XO,OO_1,XO_1$ are perpendicular to $AB,AC,$ and their bisector, $XOO_1$ is isosceles with $XO=OO_1$, and $\angle XOO_1=180-\angle A$. Also
 \[\angle AOY=2\angle ACY=2(90-\angle PAC)=2(90-\frac{A}{2})=180-\angle A = \angle XOO_1\]But $YO=OA$ as well, and $\angle YOX=\angle AOO_1$, so $\triangle OYX\cong \triangle OAO_1$. Thus $XY=AO_1=AO$.

 



\textbf{Solution 2}

Since $AHPC$ is a cyclic quadrilateral, $\angle AHC = \angle APC$. $\angle AHC = 90^\circ + \angle ABC$ and $\angle APC = 90^\circ + \angle AYC$, we find $\angle ABC = \angle AYC$. That is, $ABYC$ is a cyclic quadrilateral. Let $D$ be mid-point of $\overline{AB}$. $O, X, D$ are collinear and $OX \perp AB$. Let $M$ be second intersection of $AP$ with circumcircle of the triangle $ABC$. Let $YP \cap AC = E$, $YM \cap AB = F$. Since $M$ is mid-point of the arc $BC$, $OM\perp BC$. Since $AYMC$ is a cyclic quadrilateral, $\angle CYM = \angle CAM = \angle BAC /2$. Since $Y$ is the orthocenter of triangle $APC$, $\angle PYC = \angle CAP = \angle BAC /2$. Thus, $\angle PYM = \angle BAC$ and $AEYF$ is a cyclic quadrilateral. So, $YF \perp AB$ and $OX \parallel MY$. We will prove that $XYMO$ is a parallelogram.

 \href{https://wiki-images.artofproblemsolving.com//7/7b/Usamo2014-5.png}{https://wiki-images.artofproblemsolving.com//7/7b/Usamo2014-5.png} (figure link)

 We see that $YPM$ is an isosceles triangle and $YM=YP$. Also $XB=XP$ and $\angle BXP = 2\angle BAP = \angle BAC = \angle PYM$. Then, $BXP \sim MYP$. By spiral similarity, $BPM \sim XPY$ and $\angle XYP = \angle BMP = \angle BCA$. Hence, $\angle XYP = \angle BCA$, $XY \perp BC$. Since $OM \perp BC$, we get $XYMO$ is a parallelogram. As a result, $OM = XY$.

 (Lokman GÖKÇE)

 


	\item (\textit{USAMO 2014}) Prove that there is a constant $c>0$ with the following property: If $a, b, n$ are positive integers such that $\gcd(a+i, b+j)>1$ for all $i, j\in\{0, 1, \ldots n\}$, then
 \[\min\{a, b\}>c^n\cdot n^{\frac{n}{2}}.\]


 



\textbf{Solution}

The following solution is due to Gabriel Dospinescu and v\_Enhance (also known as Evan Chen).

 Let $N = n+1$ and assume $N$ is (very) large. We construct an $N \times N$ with cells $(i,j)$ where $0 \le i, j \le n$ and in each cell place a prime $p$ dividing $\gcd (a+i, b+j)$.

 The central claim is at least $50\%$ of the primes in this table exceed $0.001n^2$. We count the maximum number of squares they could occupy:
 \[\sum_p \left\lceil \frac{N}{p} \right\rceil^2  		\le \sum_p \left( \frac Np + 1 \right)^2 \\ 		= N^2 \sum_p \frac{1}{p^2} + 2N \sum_p \frac1p + \sum_p 1.\] Here the summation runs over primes $p \le 0.001n^2$.

 Let $r = \pi(0.001n^2)$ denote the number of such primes. Now we apply the three bounds: 
 \[\sum_p \frac{1}{p^2} < \frac 12\] which follows by adding all the primes directly with some computation, 
 \[\sum_p \frac 1p < \sum_{k=1}^r \frac 1k = O(\ln r) < o(N)\] using the harmonic series bound, and 
 \[\sum_p 1 < r \sim O \left( \frac{N^2}{\ln N} \right) < o(N^2)\] via Prime Number Theorem.  Hence the sum in question is certainly less than $\frac 12 N^2$ for $N$ large enough, establishing the central claim.

 Hence some column $a+i$ has at least one half of its primes greater than $0.001n^2$. Because this is greater than $n$ for large $n$, these primes must all be distinct, so $a+i$ exceeds their product, which is larger than 
 \[\left( 0.001n^2 \right)^{N/2} > c^n \cdot n^n\] where $c$ is some constant.

 


\end{enumerate}